\section{Christian Initiation in Clementine Literature}

This is a continuation\footnote{\url{https://gornahoor.net/?p=4636}} of thoughts from my source\footnote{\url{http://www.scribd.com/doc/13561289/-The-Authentic-Peter-The-Preaching-of-Simeon-Kefa-from-the-Journal-of-T-Flavius-Clemens-Clement-The-Clementine-Recognitions-Recognitions-of-Cl}}, as well as a running summary/commentary on the events described. \emph{Please note the tremendous import of one claim made in this treatise: Baptism is a powerful rite because \textbf{Christ the Only-Begotten made the waters} (upper \& lower)}. Peter is initiating Clement (who can be viewed as the alter-ego of Simon Magus: Simon is actually more worthy from the world's point of view, as Clement has only dreamed of attempting what Magus has already done -- Clement repents, \& follows Peter ``because he has the words of life'' -- this leads Peter to explain the lesser mysteries to Clement -- a debate with Simon Magus is the centerpiece). As we've seen, both sides have eloquent rhetoric -- Magus is more crafty and dangerous in debate, Peter more solid and sensible, but equally keen.

Simon has cast doubt on Jesus' veracity by questioning whether He contradicts Himself, and Peter is explaining what ``division'' of a house, or by the Divine sword of Truth, means in a religion based ostensibly on ``uniting the house'' or gathering in all the tribes of Israel and nations of men:

\begin{quotex}
But if any one say, How does it seem right for men to be separated from their parents? I will tell you how. Because, if they remained with them in error, they would do no good to them, and they would themselves perish with them. It is therefore right, and very right, that he who will be saved be separated from him who will not. But observe this also, that this separation does not come from those who understand aright; for they wish to be with their relatives, to do them good, and to teach them better things. But it is the vice peculiar to ignorance that it will not bear to have near it the light of truth, which confutes it; and therefore that separation originates with them. For those who receive the knowledge of the truth, because it is full of goodness, desire, if it be possible, to share it with all, as given by YHWH; yea, even with those who hate and persecute them: for they know that ignorance is the cause of their sin.
\end{quotex}

Peter thus deals with Christian paradox: the Scriptures ``speak'' in parables after the fashion of men at times, at other times not -- only a spiritual man can reconcile seemingly contradictory passages. This point has been made many times by the Patristic fathers and other theologians. Gnosticism is based on an over-simplification, or over literalization of Scripture. For instance, it is not necessary to posit a Demi-Urge Creator God opposed to the Light -- one could (for instance) simply accept the old idea of (Mouravieff's) General Law and the Hierarchies, who tutor men until they are ready. Thus, we might quote Joubert\footnote{\url{https://en.wikipedia.org/wiki/Joseph\_Joubert}}: ``Force and Right rule the world; Force, till Right is ready''. There is no necessary contradiction, but paradox. Gnosticism is thus both intellectually oversimplified, and imaginatively impoverished, while also becoming more cumbersome and complex in these areas to compensate. This is not to deny Gnostic \emph{insights}, but to reject Gnostic theology.

In a tantalizing passage, Peter holds out the idea that all religions are potentially salvific, thus agreeing with Augustine\footnote{\url{http://www.scribd.com/doc/86329463/Augustine-on-True-Religion-De-Vera-Religione}}:

\begin{quotex}
it so: you cannot know what Elohim is, but you can very easily know what Elohim is not. For how can any man fail to know that wood, or stone, or brass, or other such matter, is not Elohim? But if you will not give your mind to consider the things which you might easily apprehend, it is certain that you are hindered in the knowledge of~Elohim, not by impossibility, but by indolence; for if you had wished it, even from these useless images you might have been set on the way of understanding.
\end{quotex}


One can't help but think of the Three Trials\footnote{\url{https://gornahoor.net/?p=1902}} here: Man has to learn to say ``this is not me'' before he is worthy of Light. In the Christian tradition, there is Purification, lllumination, then Theosis.

The reason why Baptism is a rite, not a ``mere'' sign, but must be accompanied by initiation/confirmation:

\begin{quotex}
But water was made at first by the Only-begotten... His will: that you be born anew by means of waters, which were the first created... \emph{Wherefore,with prepared minds}, approach as sons to a father, that your sins may be washed away, and it may be proven before Elohim that ignorance was their sole cause...
\end{quotex}


Then follow some Jewish teachings on Torah which have some basis in Tradition: rules of chastity, etc. A very good point is made about being ``reactive'' towards other religions or points of view:

\begin{quotex}
Or shall we be so foolish, that if we see a~ worshipper of idols to be sober, we shall refuse to be sober, lest we should seem to do the same things as he does who worships idols?
\end{quotex}


The source of judgement on the Jews:

\begin{quotex}
\emph{And truly confusion is our worthy portion, if we have done no more than those who are inferior to us in knowledge... }
\end{quotex}


At the end of the debate, Peter confers the rite on Clement, arranges for a local bishop, orders the Church, and baptizes others:

\begin{quotex}
at the last he ordered me to fast; and after the fast he conferred on me the \emph{mikvah} of ever-flowing water, in the fountains adjoining the sea.
\end{quotex}


At times, Peter seems almost Platonic:

Or do you not know that~friends are always together, and are joined in memory, though they be separated bodily; as, on the other hand, some persons are near to one another in body, but are separate in mind?

Again,

\begin{quotex}
If any man who is a worshiper of Elohim had endured what this man's father has endured, immediately men would assign his belief as the cause of his calamities; but when these things come upon miserable Goyim, they charge their misfortunes upon fate. I call them miserable, because they are both vexed with errors here, and are deprived of future hope; whereas, when the worshippers of YHWH suffer these things, their patient endurance of them contributes to their cleansing from sin... \end{quotex}


Providence and Fate, again. Or, the Three Trials.

Then follows a minor miracle, as Peter and Clement go with others (on their tour of the region) to see some vinewood columns and statues of Phidias. Peter is more interested in an old beggar woman, \& when he inquires about her state, it transpires that this woman is Clement's long-missing mother; Clement also finds his missing two brothers (who are already in Peter's company as Niceta and Aquila).

Even more Platonism (and emphasis on chastity):

\begin{quotex}
Let us not shrink, then, from suffering along with her, for it is a sin to transgress any commandment. But let us teach our bodily senses, which are our outer senses, to be in subjection to our inner senses; and not compel our inner senses,which savor the things that are of Elohim, to follow the outer senses,which savor the things that are of the flesh.
\end{quotex}


Baptism is again emphasized:

\begin{quotex}
Therefore, as the order and reason of the mystery demanded, on the following day she was immersed in the sea, and returning to the lodging, was initiated in all the mysteries of truth in their order... \end{quotex}


Peter is always kind, compassionate, and ``ecumenical'' in his approach to other people (but there are specific, doctrinal and spiritual reasons for this):

\begin{quotex}
Since you seem to me to be a learned and compassionate man, inasmuch as you have come to us and wish that to be known to us which you consider to be good, we also wish to expound to you what things we believe to be good and right. And if you do not think them true, you will take in good part our good intentions towards you as we do yours towards us.
\end{quotex}


Anyone sincerely hungering for Truth is talked to by Peter, kindly -- the hunger is the thing, not the outward appearance. It is also interesting from this comment that Peter viewed ``avatars'' in a complex light (which should be clear from earlier comments):

\begin{quotex}
Those who speak the word of truth, and who enlighten the spirits of men, seem to me to be like the rays of the sun, which, when once they have come forth and appeared to the world, can no longer be concealed or hidden, while they are not so much seen by men, as they afford sight to all. Therefore it was well said by One of the heralds of the truth, Ye are the light of the world,and a city set upon a hill cannot be hid; neither do men light a candle and put it under a bushel, but upon a candlestick, that it may enlighten all who are in the house.
\end{quotex}


A new debate with an old hermit commences ‐ the old hermit maintains that the gods or God does not rule the world, but only Chance and something he calls ``Genesis''. We will continue this next week... What can be learned from these sections? Again we see the influence of Greek thought and learning on ``Torah'', almost to the point of verbatim quotes from Plato, as well as certain practices like fasting, chastity, and the performance of rites. We also see that Peter holds to the ``many truths'' position, \emph{so long as} the one truth holds paramount, and so long as there is a personal hunger for truth.


\flrightit{Posted on 2012-08-17 by Logres}

